\section{Water Orientation}
The orientation of water within the aqueous/organic interface of the system was defined using the angles formed by molecular axes and the fixed long-axis of the system (perpendicular to the interfacial plane). The two molecular axes used are the water bisector, pointing from the hydrogen side to the oxygen, and the vector normal to the plane formed by the three atoms of the water molecule. Figure \ref{fig:water-angles} illustrates the angles of interest relative to the system reference axis (i.e. Z-axis) as the molecule takes on various orientations.

% find references of the orientation of water near salt-interfaces for the above paragraph

\begin{figure}[h!]
\begin{center}
	\includegraphics[scale=0.26]{images/h2o-2dhistograms.png}
	\caption{Orientation profiles of interfacial water molecules. Two profiles were developed for characterizing interfacial water orientation based on the bisector and molecular normal vector angles formed with the system references axis (see Figure \ref{fig:water-angles}). Each profile shows the water population for each cosine of the respective angle as a function of distance to the Gibbs dividing surface (located at 0.0). The profiles in the left column are of the bisector vector, and those on the right are of the molecular normal vector. Neat-H$_2$O, NaCl, NaNO$_3$, and Na$_2$SO$_4$ system profiles are ordered from the top to bottom row, respectively.}
	\label{fig:system-orientations}
\end{center}
\end{figure}


